% =============================================================================
% Practico 01 - Maquinas de Turing
% Computabilidad y Complejidad (Teoria de la Computacion II)
% FCEFQyN - Universidad Nacional de Rio Cuarto
% Primer Cuatrimestre de 2026
% -----------------------------------------------------------------------------
% Compilacion sugerida: pdflatex practico-01.tex
% =============================================================================

\documentclass[11pt,a4paper]{article}

% --- Codificacion y idioma --------------------------------------------------
\usepackage[utf8]{inputenc}
\usepackage[T1]{fontenc}
\usepackage[spanish,es-tabla]{babel}
\usepackage{lmodern}
\usepackage{microtype}

% --- Matematica -------------------------------------------------------------
\usepackage{amsmath,amssymb,amsthm,mathtools}
\usepackage{stmaryrd}

% --- Layout y tipografia ----------------------------------------------------
\usepackage[margin=2.7cm]{geometry}
\usepackage{enumitem}
\setlist[itemize]{topsep=2pt,itemsep=2pt,parsep=0pt}
\setlist[enumerate]{topsep=2pt,itemsep=2pt,parsep=0pt}
\usepackage{parskip}
\usepackage{xcolor}
\usepackage{booktabs}
\usepackage{array}

% --- Hyperref ---------------------------------------------------------------
\usepackage[hidelinks]{hyperref}
\hypersetup{
  pdftitle={Practico 01 - Maquinas de Turing},
  pdfauthor={Computabilidad y Complejidad - UNRC 2026},
  pdfsubject={Resolucion del Practico 01},
  pdfkeywords={maquinas de Turing, decidibilidad, reconocibilidad, Sipser}
}

% --- Ambientes de teoremas --------------------------------------------------
\theoremstyle{plain}
\newtheorem{teorema}{Teorema}[section]
\newtheorem{proposicion}[teorema]{Proposicion}
\newtheorem{lema}[teorema]{Lema}
\newtheorem{corolario}[teorema]{Corolario}

\theoremstyle{definition}
\newtheorem{definicion}[teorema]{Definicion}
\newtheorem{ejemplo}[teorema]{Ejemplo}

\theoremstyle{remark}
\newtheorem*{observacion}{Observacion}
\newtheorem*{idea}{Idea}

% --- Comandos de notacion ---------------------------------------------------
\newcommand{\MT}{\mathrm{MT}}
\newcommand{\TM}{\mathrm{TM}}
\newcommand{\enc}[1]{\langle #1\rangle}
\newcommand{\lang}{\mathcal{L}}
\newcommand{\Sstar}{\Sigma^{*}}
\newcommand{\qaccept}{q_{\text{accept}}}
\newcommand{\qreject}{q_{\text{reject}}}
\newcommand{\qseek}{q_{\text{seek}}}
\newcommand{\qfindone}{q_{\text{find1}}}
\newcommand{\qfindzero}{q_{\text{find0}}}
\newcommand{\qback}{q_{\text{back}}}
\newcommand{\qscanone}{q_{\text{scan1}}}
\newcommand{\qbackone}{q_{\text{back1}}}
\newcommand{\qfindzeroa}{q_{\text{find0a}}}
\newcommand{\qfindzerob}{q_{\text{find0b}}}
\newcommand{\qbackf}{q_{\text{backf}}}
\newcommand{\qcheckzero}{q_{\text{check0}}}
\newcommand{\TMstd}{\TM_{\text{std}}}
\newcommand{\TMtwo}{\TM_{2}}
\newcommand{\TMone}{\TM_{1}}
\newcommand{\TMk}{\TM_{k}}
\newcommand{\countz}{\#_{0}}
\newcommand{\countu}{\#_{1}}
\newcommand{\blank}{\sqcup}

% --- Cabecera ---------------------------------------------------------------
\title{\textbf{Practico 01 -- Maquinas de Turing}\\[0.3em]
       \normalsize Computabilidad y Complejidad (Teor\'ia de la Computaci\'on II)\\
       \normalsize FCEFQyN -- Universidad Nacional de R\'io Cuarto\\
       \normalsize Primer Cuatrimestre de 2026}
\author{}
\date{}

\begin{document}
\maketitle
\thispagestyle{empty}

\hrule
\vspace{0.4em}
\begin{center}
\small
Resolucion completa del Practico 01. Apoyada en el Cap\'itulo 3 del libro de
M. Sipser, \emph{Introduction to the Theory of Computation}, y en el material
te\'orico de la c\'atedra (\texttt{Diag.pdf}, \texttt{MaqTuring.pdf}).
\end{center}
\vspace{0.4em}
\hrule

\tableofcontents
\bigskip
\hrule
\bigskip

% =============================================================================
\section*{Introducci\'on}
\addcontentsline{toc}{section}{Introducci\'on}

El presente trabajo resuelve la totalidad de los ejercicios del Pr\'actico 01
de la asignatura, dedicado al modelo de M\'aquinas de Turing (MT), sus
variantes y los conceptos b\'asicos de \emph{decidibilidad} y
\emph{reconocibilidad}. La gu\'ia te\'orica utilizada es:

\begin{itemize}
  \item el Cap\'itulo 3 del libro de Sipser (\texttt{chapter-03.pdf}), que
        introduce la definici\'on formal de MT (Definici\'on 3.3), sus
        ejemplos cl\'asicos ($M_1, M_2, M_3, M_4$), las variantes (multicinta,
        no determinista, doblemente infinita) y el resultado de robustez
        (Teoremas 3.13, 3.16);
  \item las l\'aminas de la c\'atedra (\texttt{Diag.pdf}), que tratan
        cardinalidad, conjuntos contables, diagonalizaci\'on de Cantor,
        existencia de lenguajes no computables y el problema de la
        terminaci\'on.
\end{itemize}

\paragraph{Convenciones de notaci\'on.}

\begin{itemize}
  \item $\Sigma$ es el alfabeto de entrada y $\Gamma\supseteq\Sigma\cup\{\blank\}$
        el alfabeto de cinta. El s\'imbolo $\blank$ denota el blanco.
  \item Una \emph{configuraci\'on} se escribe $u\,q\,v$, con $u,v\in\Gamma^*$
        y $q\in Q$, e indica que el contenido de la cinta es $uv$ (seguido
        de blancos), el estado actual es $q$ y el cabezal est\'a sobre el
        primer s\'imbolo de $v$.
  \item $L(M)$ es el lenguaje reconocido por $M$.
  \item ``Decidible'' = Turing-decidible (alguna MT lo decide); ``reconocible''
        = Turing-reconocible (alguna MT lo reconoce).
\end{itemize}

\bigskip\hrule\bigskip

% =============================================================================
\section{Ejercicio 1: lectura del Cap\'itulo 3}
\label{sec:ej1}

El primer \'item del pr\'actico consiste en \emph{leer el Cap\'itulo 3 del
libro de Sipser}. Se asume realizado como prerrequisito te\'orico para el
resto de los ejercicios. Las nociones que se utilizan a partir del
Ejercicio~2 incluyen:

\begin{itemize}
  \item \textbf{Definici\'on formal (Sipser 3.3):} una MT es una 7-tupla
        \[
          M = (Q, \Sigma, \Gamma, \delta, q_0, \qaccept, \qreject),
        \]
        con $Q$ conjunto finito de estados, $\Sigma$ alfabeto de entrada
        (sin el blanco), $\Gamma\supseteq \Sigma\cup\{\blank\}$ alfabeto de
        cinta y funci\'on de transici\'on $\delta:Q\times\Gamma\to Q\times\Gamma\times\{L,R\}$.
  \item \textbf{Reconocibilidad y decidibilidad} (Definiciones 3.5--3.6).
  \item \textbf{Variantes (Secci\'on 3.2):} maquinas multicinta, no
        deterministas, doblemente infinitas, todas equivalentes a la MT
        est\'andar (Teoremas 3.13, 3.16, Problema 3.11).
  \item \textbf{Enumeradores (Secci\'on 3.3, Teorema 3.21):} $L$ es
        reconocible si y s\'olo si existe un enumerador que enumera $L$.
  \item \textbf{Tesis de Church-Turing}: la igualdad informal entre
        ``algoritmo'' y ``MT'' justifica describir m\'aquinas a nivel
        implementaci\'on o de alto nivel.
\end{itemize}

Cada ejercicio que sigue cita expl\'icitamente las definiciones y resultados
del cap\'itulo que invoca.

\bigskip\hrule\bigskip

% =============================================================================
\section{Ejercicio 2: configuraciones de $M_2$}
\label{sec:ej2}

\subsection*{Enunciado}

Para la m\'aquina de Turing $M_2$ dada en clases (Ejemplo 3.7 de Sipser),
dar la secuencia de configuraciones que recorre con las entradas:
(a) \texttt{0}, (b) \texttt{00}, (c) \texttt{000}, (d) \texttt{000000}.

\subsection*{La m\'aquina $M_2$}

$M_2$ es la MT que decide
\[
  A = \{0^{2^n}\mid n\ge 0\},
\]
es decir, las cadenas de ceros cuya longitud es una potencia de $2$. Su
descripci\'on de alto nivel es:

\begin{quote}
\(M_2=\) ``Con entrada $w$:
\begin{enumerate}
  \item Recorrer la cinta de izquierda a derecha tachando cada segundo $0$.
  \item Si en la etapa~1 qued\'o un solo $0$, aceptar.
  \item Si en la etapa~1 quedaron m\'as de un $0$ y la cantidad fue impar,
        rechazar.
  \item Volver al inicio de la cinta.
  \item Volver a la etapa~1.''
\end{enumerate}
\end{quote}

Cada iteraci\'on divide a la mitad la cantidad de $0$s; si la longitud no es
potencia de $2$, en alguna iteraci\'on la cantidad ser\'a impar y $M_2$
rechazar\'a; si es potencia de $2$, terminar\'a con un \'unico cero y
aceptar\'a.

\subsubsection*{Definici\'on formal}

\[
  M_2 = (Q, \Sigma, \Gamma, \delta, q_1, \qaccept, \qreject),
\]
con $Q=\{q_1, q_2, q_3, q_4, q_5, \qaccept, \qreject\}$, $\Sigma=\{0\}$,
$\Gamma=\{0, x, \blank\}$, y $\delta$ dada por (Figura~3.8 de Sipser):

\begin{center}
\begin{tabular}{c|c|c|c}
\toprule
Estado & $0$ & $x$ & $\blank$ \\
\midrule
$q_1$ & $(q_2,\blank,R)$ & $(\qreject,x,R)$    & $(\qreject,\blank,R)$ \\
$q_2$ & $(q_3,x,R)$      & $(q_2,x,R)$         & $(\qaccept,\blank,R)$ \\
$q_3$ & $(q_4,0,R)$      & $(q_3,x,R)$         & $(q_5,\blank,L)$      \\
$q_4$ & $(q_3,x,R)$      & $(q_4,x,R)$         & $(\qreject,\blank,R)$ \\
$q_5$ & $(q_5,0,L)$      & $(q_5,x,L)$         & $(q_2,\blank,R)$      \\
\bottomrule
\end{tabular}
\end{center}

Intuitivamente, $q_1$ borra el primer $0$ (lo reemplaza por blanco, que
funciona como sentinela del extremo izquierdo); $q_2$, $q_3$ y $q_4$ se
alternan para tachar uno de cada dos ceros; $q_5$ vuelve al extremo izquierdo
para reiniciar el barrido.

\subsection*{Convenci\'on para escribir configuraciones}

Usamos la convenci\'on $u\,q\,v$ de Sipser: el cabezal est\'a sobre el primer
s\'imbolo de $v$ y el contenido completo de la cinta es $uv$ seguido de
blancos.

\subsection*{(a) Entrada \texttt{0}}

\begin{align*}
C_0 &= q_1\,0 \\
C_1 &= \blank\,q_2\,\blank
       &&\delta(q_1,0)=(q_2,\blank,R)\\
C_2 &= \blank\blank\,\qaccept\,\blank
       &&\delta(q_2,\blank)=(\qaccept,\blank,R)
\end{align*}

\textbf{Resultado:} acepta. Como $|0|=1=2^0$, la entrada pertenece a $A$.

\subsection*{(b) Entrada \texttt{00}}

\begin{align*}
C_0 &= q_1\,00 \\
C_1 &= \blank\,q_2\,0           &&\delta(q_1,0)=(q_2,\blank,R)\\
C_2 &= \blank x\,q_3\,\blank    &&\delta(q_2,0)=(q_3,x,R)\\
C_3 &= \blank\,q_5\,x           &&\delta(q_3,\blank)=(q_5,\blank,L)\\
C_4 &= q_5\,\blank x            &&\delta(q_5,x)=(q_5,x,L)\\
C_5 &= \blank\,q_2\,x           &&\delta(q_5,\blank)=(q_2,\blank,R)\\
C_6 &= \blank x\,q_2\,\blank    &&\delta(q_2,x)=(q_2,x,R)\\
C_7 &= \blank x\blank\,\qaccept\,\blank &&\delta(q_2,\blank)=(\qaccept,\blank,R)
\end{align*}

\textbf{Resultado:} acepta. Como $|00|=2=2^1$, pertenece a $A$.

\subsection*{(c) Entrada \texttt{000}}

\begin{align*}
C_0 &= q_1\,000 \\
C_1 &= \blank\,q_2\,00          &&\delta(q_1,0)=(q_2,\blank,R)\\
C_2 &= \blank x\,q_3\,0         &&\delta(q_2,0)=(q_3,x,R)\\
C_3 &= \blank x0\,q_4\,\blank   &&\delta(q_3,0)=(q_4,0,R)\\
C_4 &= \blank x0\blank\,\qreject\,\blank &&\delta(q_4,\blank)=(\qreject,\blank,R)
\end{align*}

\textbf{Resultado:} rechaza. La cantidad de $0$s es $3$, impar y mayor que
$1$; $M_2$ detecta esto al estar en $q_4$ y leer un blanco (``fin de pasada
con cantidad impar'').

\subsection*{(d) Entrada \texttt{000000}}

Esta entrada genera una traza m\'as larga porque $M_2$ realiza dos pasadas
completas antes de rechazar. La primera pasada deja tres ceros a procesar;
la segunda detecta que $3$ es impar.

\begin{align*}
C_0 &= q_1\,000000 \\
C_1 &= \blank\,q_2\,00000        &&\delta(q_1,0)=(q_2,\blank,R)\\
C_2 &= \blank x\,q_3\,0000       &&\delta(q_2,0)=(q_3,x,R)\\
C_3 &= \blank x0\,q_4\,000       &&\delta(q_3,0)=(q_4,0,R)\\
C_4 &= \blank x0x\,q_3\,00       &&\delta(q_4,0)=(q_3,x,R)\\
C_5 &= \blank x0x0\,q_4\,0       &&\delta(q_3,0)=(q_4,0,R)\\
C_6 &= \blank x0x0x\,q_3\,\blank &&\delta(q_4,0)=(q_3,x,R)\\
C_7 &= \blank x0x0\,q_5\,x       &&\delta(q_3,\blank)=(q_5,\blank,L)\\
C_8 &= \blank x0x\,q_5\,0x       &&\delta(q_5,x)=(q_5,x,L)\\
C_9 &= \blank x0\,q_5\,x0x       &&\delta(q_5,0)=(q_5,0,L)\\
C_{10} &= \blank x\,q_5\,0x0x    &&\delta(q_5,x)=(q_5,x,L)\\
C_{11} &= \blank\,q_5\,x0x0x     &&\delta(q_5,0)=(q_5,0,L)\\
C_{12} &= q_5\,\blank x0x0x      &&\delta(q_5,x)=(q_5,x,L)\\
C_{13} &= \blank\,q_2\,x0x0x     &&\delta(q_5,\blank)=(q_2,\blank,R)\\
C_{14} &= \blank x\,q_2\,0x0x    &&\delta(q_2,x)=(q_2,x,R)\\
C_{15} &= \blank xx\,q_3\,x0x    &&\delta(q_2,0)=(q_3,x,R)\\
C_{16} &= \blank xxx\,q_3\,0x    &&\delta(q_3,x)=(q_3,x,R)\\
C_{17} &= \blank xxx0\,q_4\,x    &&\delta(q_3,0)=(q_4,0,R)\\
C_{18} &= \blank xxx0x\,q_4\,\blank &&\delta(q_4,x)=(q_4,x,R)\\
C_{19} &= \blank xxx0x\blank\,\qreject\,\blank &&\delta(q_4,\blank)=(\qreject,\blank,R)
\end{align*}

\textbf{Resultado:} rechaza. Tras la primera pasada quedan tres ceros sin
tachar; en la segunda pasada el patr\'on de paridad cambia y $M_2$ detecta
la cantidad impar al llegar al blanco final estando en $q_4$.

\subsection*{Verificaci\'on}

Como $1=2^0$ y $2=2^1$ son potencias de $2$ pero $3$ y $6$ no lo son,
\[
  0\in A,\quad 00\in A,\quad 000\notin A,\quad 000000\notin A,
\]
en concordancia con los resultados obtenidos.

\bigskip\hrule\bigskip

% =============================================================================
\section{Ejercicio 3: deciders por marcado y emparejamiento}
\label{sec:ej3}

\subsection*{Enunciado}

Dar definiciones de m\'aquinas de Turing para decidir los siguientes
lenguajes:
\begin{enumerate}
  \item[(a)] $L_1 = \{w\in\{0,1\}^* \mid \countz(w)=\countu(w)\}$.
  \item[(b)] $L_2 = \{w\in\{0,1\}^* \mid \countz(w) = 2\cdot \countu(w)\}$.
\end{enumerate}

\subsection*{Estrategia general}

Usamos la t\'ecnica de \emph{marcado y emparejamiento} ilustrada en los
Ejemplos~3.9 y 3.11 de Sipser. El alfabeto de cinta es
$\Gamma=\{0,1,x,y,\blank\}$, donde $x$ marca un $0$ ya emparejado, $y$ marca
un $1$ ya emparejado, y $\blank$ es el blanco. Cada iteraci\'on toma un
s\'imbolo no marcado y busca su contraparte (uno o dos $0$s, seg\'un el
caso).

\subsection*{(a) MT que decide $L_1$}

\subsubsection*{Descripci\'on a nivel implementaci\'on}

\(M_{eq}=\) ``Con entrada $w$:
\begin{enumerate}
  \item Volver al inicio de la cinta.
  \item Recorrer la cinta hacia la derecha hasta el primer s\'imbolo no
        marcado ($0$ o $1$).
  \item Si no hay ninguno (s\'olo $x$, $y$ y blancos), \emph{aceptar}.
  \item Si el s\'imbolo es $0$, marcarlo como $x$ y buscar a su derecha
        el primer $1$ no marcado. Si lo encuentra, marcarlo como $y$ y
        volver a (2). Si llega a un blanco sin encontrarlo, \emph{rechazar}.
  \item Si el s\'imbolo es $1$, marcarlo como $y$ y buscar a su derecha
        el primer $0$ no marcado. Si lo encuentra, marcarlo como $x$ y
        volver a (2). Si llega a un blanco sin encontrarlo, \emph{rechazar}.''
\end{enumerate}

\subsubsection*{Descripci\'on formal}

\[
  M_{eq} = (Q, \Sigma, \Gamma, \delta, \qseek, \qaccept, \qreject),
\]
con $\Sigma=\{0,1\}$, $\Gamma=\{0,1,x,y,\blank\}$ y
$Q = \{\qseek, \qfindone, \qfindzero, \qback, \qaccept, \qreject\}$.

\paragraph{Funci\'on de transici\'on.}

\begin{align*}
\delta(\qseek,x) &= (\qseek,x,R), &
\delta(\qseek,y) &= (\qseek,y,R),\\
\delta(\qseek,0) &= (\qfindone,x,R), &
\delta(\qseek,1) &= (\qfindzero,y,R),\\
\delta(\qseek,\blank) &= (\qaccept,\blank,R). & &
\end{align*}

\begin{align*}
\delta(\qfindone,a) &= (\qfindone,a,R) \quad \text{para } a\in\{0,x,y\},\\
\delta(\qfindone,1) &= (\qback,y,L),\quad
\delta(\qfindone,\blank) = (\qreject,\blank,R).
\end{align*}

\begin{align*}
\delta(\qfindzero,a) &= (\qfindzero,a,R) \quad \text{para } a\in\{1,x,y\},\\
\delta(\qfindzero,0) &= (\qback,x,L),\quad
\delta(\qfindzero,\blank) = (\qreject,\blank,R).
\end{align*}

\begin{align*}
\delta(\qback,a) &= (\qback,a,L) \quad \text{para } a\in\{0,1,x,y\},\\
\delta(\qback,\blank) &= (\qseek,\blank,R).
\end{align*}

\subsubsection*{Correctitud y terminaci\'on}

\textbf{Si $M_{eq}$ acepta:} en cada iteraci\'on completa se marcaron en
pareja un $0$ y un $1$. La aceptaci\'on ocurre cuando $\qseek$ lee un
blanco, indicando que ya no hay s\'imbolos sin marcar. Como las marcas son
en pareja, $\countz(w)=\countu(w)$.

\textbf{Si $\countz(w)=\countu(w)$:} en cada iteraci\'on, al elegir un
s\'imbolo no marcado, la igualdad de cantidades garantiza que existe la
contraparte. As\'i $M_{eq}$ no entra en $\qreject$ y termina aceptando al
consumir todos los s\'imbolos.

\textbf{Terminaci\'on:} cada iteraci\'on completa marca al menos dos
s\'imbolos nuevos. Como $|w|<\infty$, tras a lo sumo $\lfloor |w|/2\rfloor$
iteraciones la m\'aquina se detiene.

\subsection*{(b) MT que decide $L_2$}

\subsubsection*{Descripci\'on a nivel implementaci\'on}

\(M_{2{:}1}=\) ``Con entrada $w$:
\begin{enumerate}
  \item Volver al inicio de la cinta.
  \item Buscar el primer $1$ no marcado.
  \item Si no hay $1$ sin marcar, pasar a la fase final:
        \begin{itemize}
          \item recorrer la cinta y verificar que tampoco haya ning\'un $0$
                sin marcar;
          \item si no queda ninguno, aceptar; en caso contrario, rechazar.
        \end{itemize}
  \item Si se encontr\'o un $1$, marcarlo como $y$.
  \item Volver al inicio y buscar el primer $0$ sin marcar; si no existe,
        rechazar; si existe, marcarlo como $x$.
  \item Continuar a la derecha y buscar un segundo $0$ sin marcar; si no
        existe, rechazar; si existe, marcarlo como $x$.
  \item Volver al inicio y volver a (2).''
\end{enumerate}

\subsubsection*{Descripci\'on formal}

\[
  M_{2{:}1} = (Q, \Sigma, \Gamma, \delta, \qscanone, \qaccept, \qreject),
\]
con $\Sigma=\{0,1\}$, $\Gamma=\{0,1,x,y,\blank\}$ y
\[
  Q = \{\qscanone, \qbackone, \qfindzeroa, \qfindzerob, \qback, \qbackf,
        \qcheckzero, \qaccept, \qreject\}.
\]

\paragraph{Funci\'on de transici\'on.}

\begin{align*}
\delta(\qscanone,a) &= (\qscanone,a,R) \quad \text{para } a\in\{0,x,y\},\\
\delta(\qscanone,1) &= (\qbackone,y,L),\\
\delta(\qscanone,\blank) &= (\qbackf,\blank,L).
\end{align*}

\begin{align*}
\delta(\qbackone,a) &= (\qbackone,a,L) \quad \text{para } a\in\{0,1,x,y\},\\
\delta(\qbackone,\blank) &= (\qfindzeroa,\blank,R).
\end{align*}

\begin{align*}
\delta(\qfindzeroa,a) &= (\qfindzeroa,a,R) \quad \text{para } a\in\{1,x,y\},\\
\delta(\qfindzeroa,0) &= (\qfindzerob,x,R),\quad
\delta(\qfindzeroa,\blank) = (\qreject,\blank,R).
\end{align*}

\begin{align*}
\delta(\qfindzerob,a) &= (\qfindzerob,a,R) \quad \text{para } a\in\{1,x,y\},\\
\delta(\qfindzerob,0) &= (\qback,x,L),\quad
\delta(\qfindzerob,\blank) = (\qreject,\blank,R).
\end{align*}

\begin{align*}
\delta(\qback,a) &= (\qback,a,L) \quad \text{para } a\in\{0,1,x,y\},\\
\delta(\qback,\blank) &= (\qscanone,\blank,R).
\end{align*}

\begin{align*}
\delta(\qbackf,a) &= (\qbackf,a,L) \quad \text{para } a\in\{0,1,x,y\},\\
\delta(\qbackf,\blank) &= (\qcheckzero,\blank,R).
\end{align*}

\begin{align*}
\delta(\qcheckzero,x) &= (\qcheckzero,x,R),\quad
\delta(\qcheckzero,y) = (\qcheckzero,y,R),\\
\delta(\qcheckzero,0) &= (\qreject,0,R),\quad
\delta(\qcheckzero,1) = (\qreject,1,R),\\
\delta(\qcheckzero,\blank) &= (\qaccept,\blank,R).
\end{align*}

\subsubsection*{Correctitud y terminaci\'on}

\textbf{Si $M_{2{:}1}$ acepta:} en cada ciclo completado se marc\'o
exactamente un $1$ y dos $0$s. La fase final verifica que no quedan ni $1$s
ni $0$s sueltos. Llamando $k$ al n\'umero de ciclos completados,
$\countu(w)=k$ y $\countz(w)=2k$, por lo cual
$\countz(w)=2\,\countu(w)$.

\textbf{Si $\countz(w)=2\,\countu(w)$:} al iniciar el ciclo $i+1$ restan
$\countu(w)-i$ unos y $2(\countu(w)-i)$ ceros sin marcar. Mientras
$\countu(w)-i\ge 1$, hay al menos dos $0$s para emparejar al $1$ recin
marcado, asi que el ciclo completa sin rechazar. Tras $\countu(w)$ ciclos no
quedan s\'imbolos sin marcar y la fase final acepta.

\textbf{Terminaci\'on:} cada ciclo marca un $1$ y dos $0$s, o rechaza. Como
$|w|<\infty$, tras a lo sumo $\countu(w)$ ciclos la m\'aquina llega a la
fase final.

\subsection*{Conclusi\'on}

\[
  \boxed{L_1 \text{ y } L_2 \text{ son lenguajes decidibles.}}
\]

Las MT $M_{eq}$ y $M_{2{:}1}$ son deciders explicitos. Sus implementaciones
en Python (archivo \texttt{ejercicio4\_mt\_teorico.py}) verifican
autom\'aticamente la equivalencia entre la descripci\'on formal y el
comportamiento ejecutado.

\bigskip\hrule\bigskip

% =============================================================================
\section{Ejercicio 4: cinta doblemente infinita}
\label{sec:ej4}

\subsection*{Enunciado}

Una m\'aquina de Turing con \emph{cinta doblemente infinita} es una MT cuya
cinta es infinita tanto a la izquierda como a la derecha. Mostrar que estas
m\'aquinas reconocen los mismos lenguajes que las MT est\'andar
(Problema 3.11 de Sipser).

\subsection*{Modelos}

\begin{itemize}
  \item \textbf{MT est\'andar} ($\TMstd$): una sola cinta infinita hacia la
        derecha, con extremo izquierdo. Si el cabezal intenta cruzar el
        extremo izquierdo, se queda en su lugar.
  \item \textbf{MT doblemente infinita} ($\TMtwo$): una sola cinta indexada
        por $\mathbb Z$. La entrada $w=w_0\cdots w_{n-1}$ ocupa las celdas
        $0,1,\dots,n-1$, el resto contiene blancos, y el cabezal arranca en
        la celda $0$.
\end{itemize}

\begin{teorema}
\label{teo:doble}
$\lang(\TMstd) = \lang(\TMtwo)$.
\end{teorema}

\begin{proof}
Probamos las dos inclusiones por separado.

\textbf{($\subseteq$)} Sea $M\in\TMstd$. Construimos $N\in\TMtwo$ que simula
a $M$ reservando un marcador $\vdash\in\Gamma_N\setminus\Gamma_M$ en la celda
$-1$. La descripci\'on:

\begin{quote}
\(N=\) ``Con entrada $w$:
\begin{enumerate}
  \item Escribir $\vdash$ en la celda $-1$.
  \item Simular paso a paso a $M$ sobre las celdas $0,1,\dots$
  \item Si la transici\'on indica mover a la izquierda y la celda actual es
        $\vdash$, $N$ se queda en su lugar (replicando el comportamiento de
        borde izquierdo de $M$).
  \item Cuando $M$ acepta o rechaza, $N$ acepta o rechaza.''
\end{enumerate}
\end{quote}

Por simulaci\'on paso a paso, $L(N) = L(M)$.

\textbf{($\supseteq$)} Sea $B\in\TMtwo$. Construimos $S\in\TMstd$ que ``dobla
la cinta sobre s\'i misma''.

\paragraph{Codificaci\'on.}
Definimos la biyecci\'on $e:\mathbb Z\to\mathbb N$:
\[
  e(i) = \begin{cases} 2i & i\ge 0,\\ -2i-1 & i<0.\end{cases}
\]
Esto manda $0\mapsto 0,\ 1\mapsto 2,\ 2\mapsto 4,\dots,\ -1\mapsto 1,\ -2\mapsto 3,\dots$
La celda $e(i)$ de $S$ guarda el s\'imbolo que $B$ tiene en la celda $i$.
Para indicar la posici\'on del cabezal de $B$, usamos un alfabeto extendido
$\Gamma_S = \Gamma\cup\dot\Gamma$ con $\dot\Gamma=\{\dot a : a\in\Gamma\}$.
Mantenemos la invariante: en todo momento exactamente una celda de $S$ tiene
un s\'imbolo punteado, y corresponde a la posici\'on del cabezal de $B$.

\paragraph{Inicializaci\'on.}
$S$ reorganiza $w=w_0\cdots w_{n-1}$ en las celdas $0,2,\dots,2(n-1)$,
intercala blancos en las celdas $1,3,\dots$, y marca la celda $0$ como
$\dot{w_0}$.

\paragraph{Simulaci\'on de un paso.}
Si $S$ lee $\dot a$ en la celda $k=e(h)$, $B$ est\'a en estado $q$ con
$\delta(q,a)=(p,b,D)$, entonces $S$:
\begin{enumerate}
  \item escribe $b$ (sin marca) en la celda $k$;
  \item calcula $k'=e(h')$ con $h'=h\pm 1$ seg\'un $D$;
  \item marca la celda $k'$.
\end{enumerate}
Las f\'ormulas expl\'icitas para $k'$ se obtienen examinando paridades:
\[
  k'\big|_{D=R} = \begin{cases} k+2 & k\text{ par},\\ 0 & k=1,\\ k-2 & k>1\text{ impar},\end{cases}
  \qquad
  k'\big|_{D=L} = \begin{cases} 1 & k=0,\\ k-2 & k>0\text{ par},\\ k+2 & k\text{ impar}.\end{cases}
\]
$S$ s\'olo se mueve a lo sumo dos celdas para reubicar la marca.

\paragraph{Correctitud.}
Por inducci\'on en el n\'umero de pasos, despu\'es de simular el paso
$t$-\'esimo la cinta de $S$ codifica fielmente la de $B$ y la \'unica celda
marcada coincide con el cabezal de $B$. Por lo tanto, $S$ acepta sii $B$
acepta y $L(S)=L(B)$.
\end{proof}

\begin{observacion}
Este resultado es un caso m\'as de la \emph{robustez} del modelo de MT
(Sipser, p\'ag. 176): variantes ``naturales'' del modelo no cambian la clase
de lenguajes reconocidos. La codificaci\'on $e$ es la misma idea que la
biyecci\'on $\mathbb N\cong\mathbb Z$ usada en la teor\'ia de la c\'atedra
(\texttt{Diag.pdf}); aqu\'i se aplica como herramienta de simulaci\'on.
\end{observacion}

\bigskip\hrule\bigskip

% =============================================================================
\section{Ejercicio 5: subconjunto infinito decidible}
\label{sec:ej5}

\subsection*{Enunciado}

Demostrar que cualquier lenguaje infinito que es reconocible por una m\'aquina
de Turing tiene un subconjunto infinito que es decidible
(Problema 3.19 de Sipser).

\begin{proposicion}
\label{prop:sub-decidible}
Sea $A\subseteq\Sstar$ infinito y Turing-reconocible. Entonces existe
$B\subseteq A$ tal que $B$ es infinito y decidible.
\end{proposicion}

\begin{proof}
Como $A$ es Turing-reconocible, por el Teorema~3.21 de Sipser existe un
enumerador $E$ que enumera $A$:
\[
  A = \{s : E \text{ imprime } s \text{ en alg\'un momento}\}.
\]
Fijamos sobre $\Sstar$ el orden \emph{shortlex} (primero por longitud, luego
lexicogr\'afico): $u<v$ si $|u|<|v|$, o si $|u|=|v|$ y $u<_{\text{lex}}v$.
Es un orden total efectivo y bien fundado.

\paragraph{Sucesi\'on creciente en $A$.}
Definimos $(b_i)_{i\ge 1}$ inductivamente:
\begin{itemize}
  \item $b_1$ = primera palabra que imprime $E$.
  \item Dado $b_i$, sea $b_{i+1}$ la primera palabra que imprime $E$ y
        cumple $b_{i+1} > b_i$ en shortlex.
\end{itemize}
La definici\'on es efectiva porque, para cada $i$:
\begin{itemize}
  \item podemos simular $E$ paso a paso;
  \item como $A$ es infinito y $b_i\in A$, hay infinitas palabras de $A$
        mayores que $b_i$ en shortlex, y $E$ imprime cada una en tiempo
        finito;
  \item por lo tanto, en alg\'un paso $E$ imprime una mayor que $b_i$, que
        tomamos como $b_{i+1}$.
\end{itemize}
Definimos $B = \{b_1, b_2, b_3, \dots\}$. Por construcci\'on, cada $b_i\in A$,
$B\subseteq A$, y la sucesi\'on $b_1<b_2<\cdots$ es estrictamente creciente,
de modo que $B$ es infinito.

\paragraph{Decisi\'on de $B$.}
Construimos un decidor $D_B$:
\begin{quote}
\(D_B=\) ``Con entrada $x$:
\begin{enumerate}
  \item Generar $b_1, b_2, \dots$ de a uno usando el procedimiento anterior.
  \item Detenerse en el primer $b_k$ con $b_k\ge x$ en shortlex.
  \item Si $b_k=x$, aceptar; si $b_k>x$, rechazar.''
\end{enumerate}
\end{quote}

\paragraph{Correctitud.}
Si $x\in B$, existe $k$ con $b_k=x$; al alcanzar ese $k$ el algoritmo
acepta. Si $x\notin B$, sea $k$ el primer $b_k\ge x$; necesariamente
$b_k>x$, y el algoritmo rechaza.

\paragraph{Terminaci\'on.}
Bajo shortlex, $\{y\in\Sstar : y\le x\}$ es finito. Como $(b_i)$ es
estrictamente creciente, en finitos pasos se alcanza $b_k\ge x$. Cada $b_i$
se calcula en tiempo finito. Luego $D_B$ siempre se detiene.

Por lo tanto $B$ es infinito y decidible, y $B\subseteq A$.
\end{proof}

\begin{observacion}
Es interesante notar que un lenguaje reconocible no es en general decidible
(como muestra $A_{TM}$, Cap\'itulo 4 y \texttt{Diag.pdf}), pero \emph{siempre}
se puede aislar dentro de \'el una porci\'on decidible que conserve la
cardinalidad. En particular, un lenguaje admite enumeraci\'on creciente
efectiva si y s\'olo si es decidible (Problema~3.18 de Sipser).
\end{observacion}

\bigskip\hrule\bigskip

% =============================================================================
\section{Ejercicio 6: MT con input read-only}
\label{sec:ej6}

\subsection*{Enunciado}

Mostrar que una MT con una sola cinta que no puede escribir sobre la parte
de la cinta que contiene el input s\'olo puede reconocer lenguajes regulares
(Problema 3.20 de Sipser).

\subsection*{Modelo y proposici\'on}

Adoptamos la interpretaci\'on est\'andar: la MT $M$ tiene una sola cinta y,
sobre las celdas que contienen el input, las transiciones $\delta(q,a)$
\emph{deben} escribir el mismo s\'imbolo $a$. Asumimos por comodidad que la
zona del input est\'a delimitada por marcadores $\vdash$ (extremo izquierdo)
y $\dashv$ (extremo derecho), y que $M$ nunca cruza esos l\'imites.

Bajo esta hip\'otesis, la cinta es de \emph{s\'olo lectura} sobre el input,
las transiciones se reducen a $\delta(q,a)=(p,a,D)$ con $D\in\{L,R\}$, y el
modelo coincide con un \textbf{2DFA} (aut\'omata finito determinista de dos
v\'ias).

\begin{proposicion}
\label{prop:read-only}
Todo lenguaje reconocido por una MT $M$ que no escribe sobre el input es
\emph{regular}.
\end{proposicion}

\begin{proof}
Construimos un 2DFA equivalente. Sean
\[
  M = (Q, \Sigma, \Gamma, \delta, q_0, \qaccept, \qreject)
\]
y $A=(Q, \Sigma, \delta_A, q_0, \{\qaccept\})$ con $\delta_A:Q\times(\Sigma\cup\{\vdash,\dashv\})\to Q\times\{L,R\}$ definida por:

\begin{itemize}
  \item Si $\delta(q,a)=(p,a,D)$ con $a\in\Sigma$, entonces $\delta_A(q,a)=(p,D)$.
  \item En $\vdash$ y $\dashv$, $\delta_A$ reproduce el comportamiento de
        borde de $M$ (en particular, $A$ no cruza los marcadores).
  \item Cuando $M$ entra en $\qaccept$ o $\qreject$, $A$ entra en esos
        estados.
\end{itemize}

Como $M$ no altera la cinta de entrada, cada paso de $M$ se reproduce uno a
uno como un paso de $A$. Luego $L(M) = L(A)$.

Por el teorema cl\'asico de \textbf{Rabin--Scott / Shepherdson} (1959), los
lenguajes reconocidos por 2DFA son exactamente los regulares. Por lo tanto
$L(A)$ es regular y, por la igualdad, $L(M)$ tambi\'en lo es.
\end{proof}

\begin{observacion}
Este resultado ilumina por qu\'e la \emph{escritura sobre la cinta} es
esencial para la potencia de las MT. Sin esa capacidad, las MT colapsan a la
clase regular y pierden la posibilidad de reconocer lenguajes tan b\'asicos
como $\{0^n1^n : n\ge 0\}$, que requieren memoria proporcional a la entrada.
\end{observacion}

\bigskip\hrule\bigskip

% =============================================================================
\section{Parte de programaci\'on}
\label{sec:prog}

La parte de programaci\'on pide implementar las tres versiones de MT (cinta
\'unica e infinita a derecha, doblemente infinita, y multicinta), las MT del
Ejercicio 3 te\'orico, y los algoritmos de traducci\'on entre los modelos.
Las implementaciones est\'an disponibles en el repositorio
(\texttt{maquina\_turing.py}, \texttt{maquina\_turing\_doble\_infinita.py},
\texttt{maquina\_turing\_multicinta.py}, \texttt{ejercicio\{1,2,3,4\}\_*.py})
y han sido testeadas contra los predicados matem\'aticos correspondientes.

A continuaci\'on documentamos los algoritmos de traducci\'on, que constituyen
la parte conceptualmente m\'as relevante.

\subsection*{Notaci\'on}

\begin{itemize}
  \item $\TMone$: MT con una cinta infinita hacia la derecha.
  \item $\TMtwo$: MT con una cinta doblemente infinita.
  \item $\TMk$: MT con $k\ge 1$ cintas, cada una infinita hacia la derecha.
\end{itemize}

\subsection*{A) $\TMone \to \TMtwo$}

Dada $M\in\TMone$, construir $N\in\TMtwo$ con $L(N)=L(M)$.

\begin{enumerate}
  \item Reservar $\vdash\in\Gamma_N\setminus\Gamma_M$ como marcador de borde
        izquierdo.
  \item En la inicializaci\'on, $N$ escribe $\vdash$ en la celda $-1$.
  \item Cada transici\'on de $M$ se traslada literalmente a $N$. Las
        transiciones de movimiento a la izquierda sobre la celda con
        $\vdash$ dejan el cabezal en su lugar, replicando el comportamiento
        de borde izquierdo de $\TMone$.
  \item Estados de aceptaci\'on y rechazo se preservan.
\end{enumerate}

(Construcci\'on detallada en la primera parte del Ejercicio 4.)

\subsection*{B) $\TMtwo \to \TMone$}

Dada $B\in\TMtwo$, construir $S\in\TMone$ con $L(S)=L(B)$.

\begin{enumerate}
  \item Codificar $i\in\mathbb Z$ con $e(i)=2i$ si $i\ge 0$, $e(i)=-2i-1$ si
        $i<0$.
  \item En $S$, la celda $e(i)$ guarda el s\'imbolo de $B$ en la celda $i$.
  \item Marcar exactamente una celda con un s\'imbolo punteado para indicar
        la posici\'on del cabezal de $B$.
  \item Inicializaci\'on: reorganizar la entrada en las celdas $0,2,\dots,2(n-1)$,
        intercalar blancos en las impares, marcar la celda $0$.
  \item Cada paso de $B$ se simula con escritura local, c\'alculo de la
        nueva posici\'on $e(i\pm 1)$ y reubicaci\'on de la marca.
\end{enumerate}

(Construcci\'on detallada en la segunda parte del Ejercicio 4.)

\subsection*{C) $\TMone \to \TMk$}

Dada $M\in\TMone$ y $k\ge 1$, construir $N\in\TMk$ con $L(N)=L(M)$:

\begin{enumerate}
  \item Usar la cinta 1 de $N$ para simular la cinta de $M$.
  \item Mantener las cintas $2,3,\dots,k$ en blanco y sin movimiento.
  \item Cada $\delta_M(q,a)=(p,b,D)$ se traduce a una transici\'on de
        $\TMk$ que solo afecta la cinta 1.
\end{enumerate}

\subsection*{D) $\TMk \to \TMone$}

Es la construcci\'on del Teorema~3.13 de Sipser. Dada $K\in\TMk$, construir
$U\in\TMone$:

\begin{enumerate}
  \item Codificar las $k$ cintas de $K$ en una sola cinta de $U$ con
        delimitadores: $\#\,u_1\,\#\,u_2\,\#\,\cdots\,\#\,u_k\,\#$.
  \item En cada bloque, marcar con un s\'imbolo punteado el s\'imbolo bajo
        el cabezal de la cinta correspondiente.
  \item Simular un paso de $K$ en tres fases:
        \begin{itemize}
          \item \emph{Lectura:} $U$ recorre la cinta y memoriza en su control
                finito los $k$ s\'imbolos marcados.
          \item \emph{C\'alculo:} consulta $\delta_K$ con esos s\'imbolos y
                el estado actual.
          \item \emph{Escritura/movimiento:} segunda pasada para actualizar
                los s\'imbolos marcados y mover las marcas. Si una marca
                debe avanzar m\'as all\'a del fin de su bloque, $U$ expande
                el bloque desplazando todo a la derecha.
        \end{itemize}
\end{enumerate}

Por simulaci\'on paso a paso, $L(U)=L(K)$.

\subsection*{E) Traducciones restantes}

Por composici\'on:

\begin{itemize}
  \item $\TMtwo\to\TMk$: aplicar B y luego C.
  \item $\TMk\to\TMtwo$: aplicar D y luego A.
\end{itemize}

\subsection*{Conclusi\'on}

\[
  \mathcal L(\TMone) = \mathcal L(\TMtwo) = \mathcal L(\TMk)
  \quad\text{para todo } k\ge 1.
\]

Esta equivalencia es la \emph{robustez} del modelo de MT y unifica las tres
versiones implementadas en c\'odigo.

\bigskip\hrule\bigskip

% =============================================================================
\section*{Conclusiones}
\addcontentsline{toc}{section}{Conclusiones}

A lo largo del pr\'actico hemos consolidado los conceptos centrales del
Cap\'itulo~3 de Sipser:

\begin{itemize}
  \item El Ejercicio~2 muestra c\'omo opera una MT concreta paso a paso.
  \item El Ejercicio~3 ilustra la t\'ecnica de \emph{marcado y emparejamiento}
        para construir deciders de lenguajes contextuales.
  \item Los Ejercicios~4 y~6 estudian variantes del modelo: la cinta
        doblemente infinita preserva la potencia, mientras que la
        restricci\'on de no escribir sobre el input la reduce a la clase
        regular.
  \item El Ejercicio~5 establece un puente entre reconocibilidad y
        decidibilidad: dentro de cualquier lenguaje reconocible infinito hay
        un fragmento decidible infinito.
  \item La parte de programaci\'on cierra el c\'irculo conectando la
        formalizaci\'on de los modelos con su simulaci\'on en c\'odigo y los
        algoritmos de traducci\'on.
\end{itemize}

En conjunto, los resultados refuerzan la \emph{tesis de Church--Turing}: el
modelo de MT (en cualquiera de sus variantes razonables) captura
efectivamente la noci\'on intuitiva de algoritmo, y proveen los cimientos
sobre los que se construyen los conceptos de \emph{indecibilidad} y
\emph{reducci\'on} desarrollados en el Pr\'actico~02.

\end{document}
